\chapter{第八单元}

\Needspace{12\baselineskip}

\section{第八单元素养评价（A）}

\Needspace{8\baselineskip}

\pdfbookmark[2]{原书第65页}{orig-65}

\subsection*{原书第65页}

% p065_s01
\BookQuestion{一、1}
在轴对称图形下面的括号里画“√”。按从左到右、先上后下观察图中的10个图案。\par
\BookFigure{figures/p065_s01-01.tex}
\begin{BookAnswer}
能沿某条直线对折后两侧完全重合的图形是轴对称图形。按从左到右、先上后下编号，应给第1、3、5、6、7、10个图打√。第2、4是一般斜四边形，第8为波浪形，第9两臂长度不同的折箭头，都不能沿一条直线重合。
\BookFigure{figures/p065_s01-02.tex}
\end{BookAnswer}

% p065_s02
\BookQuestion{一、2}
曲阜孔庙是全国规模最大、历史最久的孔庙，1994年被列入《世界文化遗产名录》。“金声玉振”的门坊沧桑古朴，是孔庙的第一道门坊，它采用的是（\hspace{1.8em}）图形设计。\par
\BookFigure{figures/p065_s02-01.tex}
\begin{BookAnswer}
按题意讨论门坊的建筑几何结构：左右构件呈镜面对称，本题通常填“轴对称”。\par 必须说明判断范围：牌匾上的“金声玉振”四字并不因建筑结构对称而成为对称文字；若把这些文字与建筑一起作为一个完整的平面图案，则不能说整幅图关于竖直线完全重合。重绘图中的题名作为名称注释，不能据此忽略原图中有文字这一事实。
\end{BookAnswer}

% p065_s03
\BookQuestion{一、3（1）}
观察方格图。（1）“小船”先向（\hspace{1.8em}）平移了（\hspace{1.8em}）格，再向（\hspace{1.8em}）平移了（\hspace{1.8em}）格。\par
\BookFigure{figures/p065_s03-01.tex}
\begin{BookAnswer}
选择小船上同一个顶点作参照。先到下一幅船：向下5格；再到右边船：向右4格。平移只改位置，船头、船帆的朝向不变。填“下、5、右、4”。数的是同一顶点从起点到终点跨过的格数，不能数两个图形之间空了几格。若题目只限定起点和终点，也可先右移4格再下移5格；本题的四空仍按给定中间图顺序填写。
\end{BookAnswer}

% p065_s04
\BookQuestion{一、3（2）}
观察方格图。（2）“小鱼”先向（\hspace{1.8em}）平移了（\hspace{1.8em}）格，再向（\hspace{1.8em}）平移了（\hspace{1.8em}）格。\par
\BookFigure{figures/p065_s04-01.tex}
\begin{BookAnswer}
选择小鱼嘴尖作参照。先向上5格，再向左5格；所有对应点都移动相同格数。填“上、5、左、5”。
\end{BookAnswer}

% p065_s05
\BookQuestion{一、4}
我们学过的平面图形中，长方形有（\hspace{1.8em}）条对称轴，有4条对称轴的图形是（\hspace{1.8em}）形，（\hspace{1.8em}）三角形有3条对称轴。\par
\begin{BookAnswer}
长方形有2条对称轴，都是经过一组对边中点的直线；有4条对称轴的常见四边形是正方形；三角形中有3条对称轴的是等边三角形。依次填2、正方、等边。
\end{BookAnswer}

\Needspace{8\baselineskip}

\pdfbookmark[2]{原书第66页}{orig-66}

\subsection*{原书第66页}

% p066_s01
\BookQuestion{一、5}
奥运五环连在一起代表着五大洲的人们友好相处，齐聚奥运盛会。它可以由（\hspace{1.8em}）形（平移、旋转）（在正确答案下面画“√”）而成，它有（\hspace{1.8em}）条对称轴，请在图中画出来。\par
\BookFigure{figures/p066_s01-01.tex}
\begin{BookAnswer}
基本图形是圆，五个圆通过平移组成五环，应在“平移”下打√。圆本身没有方向，但图案中五个圆的圆心分布关于中间一条竖直线对称，所以组合图案有1条对称轴。将轴穿过上排中间圆的圆心与下排两圆圆心连线的中点。
\BookFigure{figures/p066_s01-02.tex}
\end{BookAnswer}

% p066_s02
\BookQuestion{一、6（1）}
找规律，画一画。（1）四格中小房子朝向依次为上、右、空白、左，在空白处补画。\par
\BookFigure{figures/p066_s02-01.tex}
\begin{BookAnswer}
小图依次顺时针旋转90°，朝向为上、右、下、左，空白格应画朝下的小房子。画之前先对应每个顶点，房子大小、形状保持不变。
\BookFigure{figures/p066_s02-02.tex}
\end{BookAnswer}

% p066_s03
\BookQuestion{一、6（2）}
找规律，画一画。（2）四格中折箭头依次旋转，在第四格空白处补画。\par
\BookFigure{figures/p066_s03-01.tex}
\begin{BookAnswer}
折箭头每次顺时针旋转90°。第四格应为竖直箭头向上、水平箭头向右的折箭头；折点和两条臂的长度都与前三幅一致。
\BookFigure{figures/p066_s03-02.tex}
\end{BookAnswer}

% p066_s04
\BookQuestion{一、7（①）}
要将图二恢复到图一，将①部分先向上（\hspace{1.8em}）（\hspace{1.8em}）格，再向（\hspace{1.8em}）（\hspace{1.8em}）（\hspace{1.8em}）格。图一、图二为4×4方格拼马图。\par
\BookFigure{figures/p066_s04-01.tex}
\begin{BookAnswer}
比较图一、图二中马下半身的格点位置：图二的①先向上平移1格，再向右平移1格，即与图一对应部分重合。填“平移、1、右、平移、1”。
\BookFigure{figures/p066_s04-02.tex}
\end{BookAnswer}

% p066_s05
\BookQuestion{一、7（②）}
要将图二恢复到图一，将②部分绕（\hspace{1.8em}）点（\hspace{1.8em}）时针旋转（\hspace{1.8em}）度，再向（\hspace{1.8em}）平移（\hspace{1.8em}）格。图二右下方方格的左下顶点为A，右上顶点为B。\par
\BookFigure{figures/p066_s05-01.tex}
\begin{BookAnswer}
固定A点，把图二的②逆时针旋转90°，使马身转到正确方向，再向上平移1格。填“A、逆、90、上、1”。旋转时A不动，平移时每个点都上移同样1格。动笔前先圈出本题的旋转中心A，再核对“逆时针、90°”；不能凭图案中心另选一个点。
\BookFigure{figures/p066_s05-02.tex}
\end{BookAnswer}

% p066_s06
\BookQuestion{二、1}
下面运动现象中，属于平移运动的有（\hspace{1.8em}）个。图示依次是抽拉键盘托架、订书机、旋转闸、门环、抽屉。A.1；B.2；C.3；D.4。\par
\BookFigure{figures/p066_s06-01.tex}
\begin{BookAnswer}
图中的键盘托架或抽拉板、抽屉作直线移动，整体方向不改变，属于平移；转闸、订书机活动部分和门环等绕固定处转动。平移有2个，选B。
\end{BookAnswer}

% p066_s07
\BookQuestion{二、2}
下面几组图形中，甲图旋转得到乙图的是（\hspace{1.8em}）。一组、二组的三角形朝向发生改变，三组的三角形朝向相同。A.一组；B.二组；C.一组和二组；D.二组和三组。\par
\BookFigure{figures/p066_s07-01.tex}
\begin{BookAnswer}
旋转后形状、大小保持不变，但朝向改变。第一组和第二组均可通过旋转让甲与乙重合；第三组的乙主要是同朝向移位。选C。
\end{BookAnswer}

% p066_s08
\BookQuestion{二、3}
左图黑白风车图案是用（\hspace{1.8em}）图形，（\hspace{1.8em}）得到的。A.空白正方形，平移；B.黑三角形，平移；C.空白三角形，旋转；D.图示黑白三角形组成的四边形拼块，旋转。\par
\BookFigure{figures/p066_s08-01.tex}
\BookFigure{figures/p066_s08-02.tex}
\begin{BookAnswer}
观察黑白拼块，以选项D的小图形为基本单元，绕公共中心旋转，能拼出风车图案；单纯平移不能改变风叶朝向。选D，第二空填“旋转”。
\end{BookAnswer}

\Needspace{8\baselineskip}

\pdfbookmark[2]{原书第67页}{orig-67}

\subsection*{原书第67页}

% p067_s01
\BookQuestion{二、4}
剪纸是我国国家级非物质文化遗产，将纸对折后剪出重复的图案是剪纸中最基础的技法，右图是下面（\hspace{1.8em}）项剪出的。选择A、B、C或D。\par
\BookFigure{figures/p067_s01-01.tex}
\begin{BookAnswer}
对折再剪时，展开图必须关于折痕完全对称。把选项图案关于折边反射，C可形成右图四个尖角且左右、上下配对的星形，选C。不能只比较某一个尖角。
\end{BookAnswer}

% p067_s02
\BookQuestion{二、5}
分针从6走到9，绕中心点按（\hspace{1.8em}）方向旋转了（\hspace{1.8em}）度。A.顺时针30；B.逆时针90；C.顺时针90；D.逆时针30。\par
\begin{BookAnswer}
钟面数字按顺时针方向排列，分针从6走到9经过3大格，每大格30°，所以转过\(30^\circ\times3=90^\circ\)。填“顺时针、90”，选C。
\end{BookAnswer}

% p067_s03
\BookQuestion{三、1（第1图）}
画出下面图形的另一半，使它成为轴对称图形。（第1图，竖直虚线为对称轴。）\par
\BookFigure{figures/p067_s03-01.tex}
\begin{BookAnswer}
按“找关键点—找对应点—连线”作图。先找给定竖直虚线，左边每个顶点沿垂直于轴的水平格线，在轴右侧找等距点；原来落在轴上的点不动。再按原边的连接顺序连点，最后检查对应点连线是否被对称轴垂直平分。
\BookFigure{figures/p067_s03-02.tex}
\end{BookAnswer}

% p067_s04
\BookQuestion{三、1（第2图）}
画出下面图形的另一半，使它成为轴对称图形。（第2图，竖直虚线为对称轴，内部小三角形也要补画。）\par
\BookFigure{figures/p067_s04-01.tex}
\begin{BookAnswer}
把外轮廓的各个拐点、内侧小三角形的三个顶点分别关于竖直虚线反射。对称点在同一水平格线上，距轴格数相等，不能只补外边而漏掉内侧小三角形。
\BookFigure{figures/p067_s04-02.tex}
\end{BookAnswer}

% p067_s05
\BookQuestion{三、1（第3图）}
画出下面图形的另一半，使它成为轴对称图形。（第3图，竖直虚线为对称轴。）\par
\BookFigure{figures/p067_s05-01.tex}
\begin{BookAnswer}
将左侧斜边、底边和内部横线的各端点逐一反射到右侧，再连接。对称轴上的端点重合；两侧斜边应左右对称，横线两侧等长。
\BookFigure{figures/p067_s05-02.tex}
\end{BookAnswer}

% p067_s06
\BookQuestion{三、2（第1图）}
画出下面图形的对称轴（有几条画几条）：第1图，三叶图案。\par
\BookFigure{figures/p067_s06-01.tex}
\begin{BookAnswer}
三片相同叶片均匀排布，每条对称轴穿过一片叶片的中线和公共中心，共3条。先画过上、下方向的轴，再画另外两条斜轴。
\BookFigure{figures/p067_s06-02.tex}
\end{BookAnswer}

% p067_s07
\BookQuestion{三、2（第2图）}
画出下面图形的对称轴（有几条画几条）：第2图，圆环与斜带组成的禁止标志。\par
\BookFigure{figures/p067_s07-01.tex}
\begin{BookAnswer}
禁止标志的圆环及斜带整体有2条对称轴：一条沿斜带的中线，另一条过圆心且与它垂直。只有外圆对称不够，还要检查斜带也完全重合。
\BookFigure{figures/p067_s07-02.tex}
\end{BookAnswer}

% p067_s08
\BookQuestion{三、2（第3图）}
画出下面图形的对称轴（有几条画几条）：第3图，扇形。\par
\BookFigure{figures/p067_s08-01.tex}
\begin{BookAnswer}
扇形的对称轴是顶点所成角的角平分线，也穿过弧的中点，共1条；不要把任意半径都当成对称轴。
\BookFigure{figures/p067_s08-02.tex}
\end{BookAnswer}

% p067_s09
\BookQuestion{三、2（第4图）}
画出下面图形的对称轴（有几条画几条）：第4图，正六边形。\par
\BookFigure{figures/p067_s09-01.tex}
\begin{BookAnswer}
正六边形共有6条对称轴：连接3组相对顶点的直线，共3条；连接3组对边中点的直线，共3条。它们均过中心。
\BookFigure{figures/p067_s09-02.tex}
\end{BookAnswer}

% p067_s10
\BookQuestion{三、3（1）}
将下面的图形绕O点顺时针旋转\(90^\circ\)画出图形，再向右平移3格画出图形。\par
\BookFigure{figures/p067_s10-01.tex}
\begin{BookAnswer}
先以O为固定点，把原图每个顶点到O的线段顺时针转90°，距离保持不变，连接得到第一幅红图。再将第一幅图每个顶点向右平移3格，连接得到右侧红图。第一步须明确旋转三要素：中心O、顺时针、90°。第二步把旋转后图形的所有点（包括与O重合的那个顶点）同向右移3格；原题标定的O仍留作定位参照。不能把它理解成再选一个新中心旋转，两步一般也不能颠倒。
\BookFigure{figures/p067_s10-02.tex}
\end{BookAnswer}

\Needspace{8\baselineskip}

\pdfbookmark[2]{原书第68页}{orig-68}

\subsection*{原书第68页}

% p068_s01
\BookQuestion{三、3（2）}
先将下面三角形绕B点逆时针旋转\(90^\circ\)后画出图形，再将长方形绕D点顺时针旋转\(90^\circ\)后画出图形，最后请画出组合图形的对称轴。\par
\BookFigure{figures/p068_s01-01.tex}
\begin{BookAnswer}
三角形绕B逆时针90°：B不动；原来在B下方的A移到B右方相同距离，C也按逆时针90°逐点确定。长方形绕D顺时针90°：D不动，上边向下转，右下顶点转到D左下方。按对应顶点连接后，组合图是“屋顶与矩形”形，竖直中心线为其对称轴。请沿格点确定每个端点，再画直线。这两次旋转各有自己的三要素：三角形以B为中心、逆时针90°；长方形以D为中心、顺时针90°。只旋转指定图形，不能把整个组合绕同一中心一起转。
\BookFigure{figures/p068_s01-02.tex}
\end{BookAnswer}

% p068_s02
\BookQuestion{三、3（3）填空}
移一移，画一画。图形①先向右平移（\hspace{1.8em}）格，再向（\hspace{1.8em}）平移（\hspace{1.8em}）格，就会与图形②拼成一个长方形。\par
\BookFigure{figures/p068_s02-01.tex}
\begin{BookAnswer}
比较小三角形与右侧凹口的对应顶点：先向右平移9格，再向上平移2格，就恰好填入凹口，拼成一个长方形。填9、“上”、2。
\BookFigure{figures/p068_s02-02.tex}
\end{BookAnswer}

% p068_s03
\BookQuestion{三、3（3）作图}
移一移，画一画。将图形①移入图形②的凹口拼成长方形，把拼成的图形绕O点逆时针旋转\(90^\circ\)，请在网格图中画出旋转后的图形。\par
\BookFigure{figures/p068_s03-01.tex}
\begin{BookAnswer}
先补上移入凹口的小三角形，确认拼成完整长方形，再将整个长方形绕O逆时针转90°。O不动，长边由水平转为竖直；各顶点到O的距离不变。结果在O的左上方。
\BookFigure{figures/p068_s03-02.tex}
\end{BookAnswer}

% p068_s04
\BookQuestion{三、4（1）}
认真观察，填一填。（1）图①是以点（\hspace{1.8em}）为中心（\hspace{1.8em}）时针旋转的，在图中相应的括号里标出点A的对应点\(A^{\prime}\)。\par
\BookFigure{figures/p068_s04-01.tex}
\begin{BookAnswer}
旋转中心为B，因为实线与虚线的共同顶点B始终不动。按弧箭头从虚线位置指向实线位置看，方向填“逆”。A的对应点要与A到B的距离相同，A′标在图①右侧括号中。\par 题图有歧义：若把标有A、C的实线三角形作为原像，虚线作为旋转后的像，则实线到虚线应为顺时针，与弧箭头的示意方向相反。应先明确采用哪一种原像约定。
\BookFigure{figures/p068_s04-02.tex}
\end{BookAnswer}

% p068_s05
\BookQuestion{三、4（2）}
认真观察，填一填。（2）图②是以点（\hspace{1.8em}）为中心（\hspace{1.8em}）时针旋转的，在图中相应的括号里标出点B的对应点\(B^{\prime}\)。\par
\BookFigure{figures/p068_s05-01.tex}
\begin{BookAnswer}
旋转中心为D，因为D不动。按弧箭头由虚线位置指向实线位置看，填“顺”。B到D的距离与B′到D的距离相等，B′标在图②右下括号中。\par 若按实线长方形为原像、虚线为像，则方向应为逆时针。两图的箭头与原像标注约定不统一，左图原空按箭头依次保留“逆、顺”，并在此说明不同读法。
\BookFigure{figures/p068_s05-02.tex}
\end{BookAnswer}

% p068_s06
\BookQuestion{四、逆操作}
将一个三角形先向右平移6格，再绕中心点顺时针旋转\(90^\circ\)，得到现在的图形。如果想让它回到原来的位置，应该怎样操作？请写完整步骤。\par
\begin{BookAnswer}
原来的操作是先向右平移6格，再顺时针旋转90°。恢复时要逆序进行：第一步绕原来同一个旋转中心逆时针旋转90°；第二步向左平移6格。若先平移再反转，旋转中心的位置关系会不同，一般回不到原处。
\end{BookAnswer}

\Needspace{12\baselineskip}

\section{第八单元素养评价（B）}

\Needspace{8\baselineskip}

\pdfbookmark[2]{原书第69页}{orig-69}

\subsection*{原书第69页}

% p069_s01
\BookQuestion{一、1}
将一个图形沿一条（\hspace{1.8em}）对折，如果两侧的部分能够完全重合，这样的图形叫作（\hspace{1.8em}），折痕所在的直线叫作它的（\hspace{1.8em}）。对应点到对称轴的距离（\hspace{1.8em}），对应点的连线与对称轴互相（\hspace{1.8em}）。\par
\begin{BookAnswer}
轴对称的定义：沿一条直线对折后，两部分完全重合，这条直线叫对称轴。依次填“直线、轴对称图形、对称轴”。对应点到轴的距离“相等”，对应点连线与轴互相“垂直”。
\end{BookAnswer}

% p069_s02
\BookQuestion{一、2}
图形平移时，图形的（\hspace{1.8em}）和（\hspace{1.8em}）不发生改变，只改变（\hspace{1.8em}）。\par
\begin{BookAnswer}
平移只改变图形的“位置”，不改变“形状”和“大小”。比较前后对应边的长度与方向，就可检验是否为平移。检验时找两组对应点：它们移动的方向和距离应相同，对应边的朝向也保持不变。
\end{BookAnswer}

% p069_s03
\BookQuestion{一、3}
下图中，是平移现象画“△”，是旋转现象画“○”。图示依次是：滑动解锁、坐摩天轮、开冰箱门、坐缆车。\par
\BookFigure{figures/p069_s03-01.tex}
\begin{BookAnswer}
滑动解锁和缆车沿轨道前进，朝向保持不变，画△；摩天轮和冰箱门绕中心或门轴转动，画○。依次为△、○、○、△。注意题目要求的是运动方式，不是图片本身的对称性。
\end{BookAnswer}

% p069_s04
\BookQuestion{一、4}
如图，李阿姨将2千克物品放在台秤上，指针从0刻度开始顺时针旋转\(90^\circ\)，再放2千克物品，指针会（\hspace{1.8em}）时针再旋转（\hspace{1.8em}）度；如果李阿姨把4千克物品全部拿走，那么指针会（\hspace{1.8em}）时针旋转（\hspace{1.8em}）度。\par
\BookFigure{figures/p069_s04-01.tex}
\begin{BookAnswer}
先放2千克顺时针90°，再放同样2千克，按图示均匀刻度再顺时针90°；拿走全部4千克，就从4千克刻度退回0，逆时针180°。填“顺、90、逆、180”。
\end{BookAnswer}

% p069_s05
\BookQuestion{一、5}
如图，若AB的长度是4厘米，则A点到对称轴的距离是（\hspace{1.8em}）厘米，点C和点F到对称轴的距离都是（\hspace{1.8em}）厘米，点D的对称点是点（\hspace{1.8em}）。\par
\BookFigure{figures/p069_s05-01.tex}
\begin{BookAnswer}
AB连接一对对称点，长度4厘米，对称轴平分AB，所以A距轴\(4\div2=2\)厘米。按方格，C、F到轴各3格；由AB跨4格等于4厘米知每格1厘米，所以都是3厘米。D的镜像为E。填2、3、E。
\end{BookAnswer}

\Needspace{8\baselineskip}

\pdfbookmark[2]{原书第70页}{orig-70}

\subsection*{原书第70页}

% p070_s01
\BookQuestion{一、6}
钟面上指针旋转一大格是（\hspace{1.8em}）度，指针从“2”到“5”是指针（\hspace{1.8em}）时针旋转（\hspace{1.8em}）度。\par
\begin{BookAnswer}
钟面一周360°，分成12大格，每大格\(360^\circ\div12=30^\circ\)。从2到5沿数字增加方向转动，是顺时针，经过3大格，共90°。填30、“顺”、90。
\end{BookAnswer}

% p070_s02
\BookQuestion{一、7}
如图，三角尺向下平移了（\hspace{1.8em}）毫米。\par
\BookFigure{figures/p070_s02-01.tex}
\begin{BookAnswer}
比较同一顶点B与B′在直尺上的位置，向下移动2厘米。题目要求毫米，另行换算：\[2\text{厘米}=(2\times10)\text{毫米}=20\text{毫米}\]原空后已有毫米，只填20。不能拿A与B′两个不对应的顶点比较。
\end{BookAnswer}

% p070_s03
\BookQuestion{二、1}
判断：三角形都有三条对称轴。（\hspace{1.8em}）\par
\begin{BookAnswer}
×。一般三角形不一定有对称轴；等腰三角形通常1条，等边三角形才有3条。反例：不等边三角形没有对称轴。
\end{BookAnswer}

% p070_s04
\BookQuestion{二、2}
判断：物体平移时，物体上的每一个点走过的路程一样长。（\hspace{1.8em}）\par
\begin{BookAnswer}
√。平移时各点沿同一方向移动相同距离，所以各点走过的路径长度相同。
\end{BookAnswer}

% p070_s05
\BookQuestion{二、3}
判断：图形经过旋转后，形状会发生改变。（\hspace{1.8em}）\par
\begin{BookAnswer}
×。旋转改变位置、朝向，形状和大小保持不变。
\end{BookAnswer}

% p070_s06
\BookQuestion{二、4}
判断：正方形的对角线是它的对称轴。（\hspace{1.8em}）\par
\begin{BookAnswer}
\(\times\)。对角线是线段，对称轴是直线。正确说法：正方形每条对角线所在的直线都是它的对称轴。
\end{BookAnswer}

% p070_s07
\BookQuestion{二、5}
判断：把图中的图形在对称轴的另一边补充完整，这图形是等腰梯形。（\hspace{1.8em}）\par
\BookFigure{figures/p070_s07-01.tex}
\begin{BookAnswer}
√。图中上下两条半底边平行，把一半关于虚线补齐后，两腰互为镜像且相等，得到等腰梯形。
\BookFigure{figures/p070_s07-02.tex}
\end{BookAnswer}

% p070_s08
\BookQuestion{三、1}
下列图形中，对称轴最多的是（\hspace{1.8em}），对称轴最少的是（\hspace{1.8em}）。A.内切而不同心的两圆；B.相交的两个等圆；C.同心的两圆；D.三角形排布的三个相等圆。\par
\BookFigure{figures/p070_s08-01.tex}
\begin{BookAnswer}
同心圆图C关于任意过圆心的直线对称，有无数条，是最多；A中一大一小相切圆只有穿过两圆圆心的一条轴，最少。填C、A。
\end{BookAnswer}

% p070_s09
\BookQuestion{三、2}
下面不是轴对称图形的是（\hspace{1.8em}）。A.长方形；B.平行四边形；C.圆；D.半圆。\par
\begin{BookAnswer}
一般平行四边形没有对称轴；长方形、圆、半圆都有轴，选B。不能把长方形、菱形等特殊情形代替题图所说的一般平行四边形。
\end{BookAnswer}

% p070_s10
\BookQuestion{三、3}
右边的轴对称图形是从纸（\hspace{1.8em}）上剪下来的。选A、B或C。\par
\BookFigure{figures/p070_s10-01.tex}
\begin{BookAnswer}
将剪纸图样沿折边反射再拼合，A的边缘轮廓可展开成题图的左右对称形状，选A。
\end{BookAnswer}

% p070_s11
\BookQuestion{三、4}
一个三角形任意一条边上的高所在的直线都是它的对称轴，这个三角形是（\hspace{1.8em}）。A.等腰三角形；B.等腰直角三角形；C.等边三角形；D.任意三角形。\par
\begin{BookAnswer}
若三角形任意一边的高所在直线都是对称轴，则三组对应边都相等，只能是等边三角形，选C。等腰直角三角形只有一条轴，不能满足“任意一边”。
\end{BookAnswer}

% p070_s12
\BookQuestion{三、5}
小东在描述一个图形：它有四条边，没有平行线，有一条对称轴，以下（\hspace{1.8em}）是小东描述的。选A、B、C或D。\par
\BookFigure{figures/p070_s12-01.tex}
\begin{BookAnswer}
A为左右对称的风筝形，有四条边，没有平行边，且有1条对称轴；B有一组平行边，C有两组，D不对称，选A。
\end{BookAnswer}

% p070_s13
\BookQuestion{三、6}
下面四幅图中，（\hspace{1.8em}）平移到位置M后能与N成轴对称图形。选A、B、C或D。\par
\BookFigure{figures/p070_s13-01.tex}
\begin{BookAnswer}
先把右边N图形关于M、N之间的竖直界线反射，所需轮廓与选项A同朝向。把A平移入M后，两边互为镜像，选A。平移不能自行翻转或转向。
\end{BookAnswer}

\Needspace{8\baselineskip}

\pdfbookmark[2]{原书第71页}{orig-71}

\subsection*{原书第71页}

% p071_s01
\BookQuestion{三、7}
下列交通标志，可以经过旋转得到的是（\hspace{1.8em}）。选A、B、C或D。\par
\BookFigure{figures/p071_s01-01.tex}
\begin{BookAnswer}
C的三个弧形箭头等间隔排布，一个箭头绕中心旋转120°后能到下一个位置，反复旋转得到整体图案，选C。
\end{BookAnswer}

% p071_s02
\BookQuestion{四、1}
下图中，点O是线段AB上的一点，将线段AB绕点O顺时针方向旋转\(90^\circ\)，旋转后的线段与原线段的位置有何关系？如果逆时针方向旋转\(90^\circ\)呢？请分别在下面的图上画出来。\par
\BookFigure{figures/p071_s02-01.tex}
\begin{BookAnswer}
顺时针旋转90°：O不动，原来在O左方的A转到O上方，B转到O下方；分别保持OA、OB的长度，标A′、B′。逆时针旋转90°时，A′在O下方，B′在O上方。两种情况所得线段都与原AB垂直。先量对应点到O的距离，再画垂线，不能凭目测把两端取成一样长。
\BookFigure{figures/p071_s02-02.tex}
\end{BookAnswer}

% p071_s03
\BookQuestion{四、2（1）}
按要求画图。（1）画出图①的另一半，使它成为一个轴对称图形。\par
\BookFigure{figures/p071_s03-01.tex}
\begin{BookAnswer}
图①以给出的水平虚线为对称轴。每个原顶点在虚线下方找等距点，再按原顺序连线。原图到轴有几格，镜像就到轴有几格；斜边的对应两端都要单独定位。对应点连线必须与水平对称轴垂直，也就是沿竖直格线找点；仅仅离轴“看着一样远”还不够。
\BookFigure{figures/p071_s03-02.tex}
\end{BookAnswer}

% p071_s04
\BookQuestion{四、2（2）}
按要求画图。（2）画出图②向下平移7格后的图形。\par
\BookFigure{figures/p071_s04-01.tex}
\begin{BookAnswer}
把图②船体、桅杆、旗帜的全部顶点向下平移7格，方向保持不变，再依次连接。船帆尖和旗杆底端必须同时移动7格，不能只移动船体。
\BookFigure{figures/p071_s04-02.tex}
\end{BookAnswer}

% p071_s05
\BookQuestion{四、2（3）}
按要求画图。（3）画出图③绕点A逆时针旋转\(90^\circ\)后的图形。\par
\BookFigure{figures/p071_s05-01.tex}
\begin{BookAnswer}
以A为旋转中心逆时针90°，A不动。原来在A左侧同水平线的点移到A下方；原来左上方的点移到A左下方。到A的距离保持不变，连接三个对应点得到红三角形。
\BookFigure{figures/p071_s05-02.tex}
\end{BookAnswer}

% p071_s06
\BookQuestion{四、3}
围绕点O，旋转三角形画一个风车。\par
\BookFigure{figures/p071_s06-01.tex}
\begin{BookAnswer}
保留原三角形，以O为中心分别再旋转90°、180°、270°，画出另外三个全等三角形。每片风叶由相同三边组成；四片共用O，且相邻朝向相差90°。先确定三个顶点后连线，避免风叶末端偏离原格点。也可从原三角形开始，连续三次绕同一O点逆时针旋转90°，补出的其余三片与顺时针画法相同。这是两种作图顺序，不是另一个风车。
\BookFigure{figures/p071_s06-02.tex}
\end{BookAnswer}

% p071_s07
\BookQuestion{四、4}
画出镜子中看到的时刻。原钟的分针指向12，时针指向4。\par
\BookFigure{figures/p071_s07-01.tex}
\begin{BookAnswer}
镜面把钟面左右反射：12仍在正上方，原来指向4的时针反射后指向8，分钟仍为0分。因此镜中为8时整，画长针向12、短针向8，两个针的共同端点落在钟心。
\BookFigure{figures/p071_s07-02.tex}
\end{BookAnswer}

\Needspace{8\baselineskip}

\pdfbookmark[2]{原书第72页}{orig-72}

\subsection*{原书第72页}

% p072_s01
\BookQuestion{五、1（图②）}
按图①的规律，补全图②中缺少的部分。图①的T形按顺时针方向逐次转\(90^\circ\)，图②已给左上Y形与右下转\(180^\circ\)后的Y形。\par
\BookFigure{figures/p072_s01-01.tex}
\begin{BookAnswer}
按图①的规律，左上图形依次顺时针旋转90°、180°、270°填入右上、右下、左下。对Y形，右上应开口朝右，左下开口朝左；转动前后的三条臂长度相同，交点落在各小格相应中心。
\BookFigure{figures/p072_s01-02.tex}
\end{BookAnswer}

% p072_s02
\BookQuestion{五、1（图③）}
按图①的规律，补全图③中缺少的部分。图③已给左上F形与右上顺时针转\(90^\circ\)后的F形。\par
\BookFigure{figures/p072_s02-01.tex}
\begin{BookAnswer}
F形按同一旋转规律补齐：右下为左上旋转180°的形状，左下为旋转270°的形状。F的长竖笔、顶横笔、中横笔分别对应，不能把镜像当作旋转。
\BookFigure{figures/p072_s02-02.tex}
\end{BookAnswer}

% p072_s03
\BookQuestion{五、2}
剪纸艺术是我国文化宝库中的优秀遗产，其技法中有对称、平移和旋转现象。判断下列三幅剪纸图案分别属于什么现象，在括号内填写。\par
\BookFigure{figures/p072_s03-01.tex}
\begin{BookAnswer}
第一幅剪纸左右两鱼互为镜像，是“对称”；第二幅动物围绕中心转动，是“旋转”；第三幅动物保持朝向重复移位，是“平移”。依次填对称、旋转、平移。
\end{BookAnswer}

% p072_s04
\BookQuestion{五、3（1）}
下面是华华画的自己家新建房屋平面图，每个小正方形的边长是2米，但图中有三处错误，你能按下面的要求帮她改正吗？（1）画出花圃的另一半，使它成为一个轴对称图形。\par
\BookFigure{figures/p072_s04-01.tex}
\begin{BookAnswer}
把花圃现有下半部关于通过中间横边的水平格线反射。四个轮廓顶点分别在上侧找等距点，再连线；完成图的上、下两半完全重合。保留原房屋网格，不另改比例。
\BookFigure{figures/p072_s04-02.tex}
\end{BookAnswer}

% p072_s05
\BookQuestion{五、3（2）}
房屋平面图中，每格边长2米。（2）将厨房先向下平移12米，再向右平移4米。\par
\BookFigure{figures/p072_s05-01.tex}
\begin{BookAnswer}
每格边长2米，向下12米就是\(12\div2=6\)格；向右4米就是\(4\div2=2\)格。先把原厨房的各顶点下移6格，再右移2格，厨房仍是原来的一格宽、两格高，在新位置标“厨房”。
\BookFigure{figures/p072_s05-02.tex}
\end{BookAnswer}

% p072_s06
\BookQuestion{五、3（3）}
房屋平面图中，每格边长2米。（3）将车库绕点C按逆时针方向旋转\(90^\circ\)。\par
\BookFigure{figures/p072_s06-01.tex}
\begin{BookAnswer}
C为原车库左下角。以C为固定点逆时针90°：原来向上的长边转为向左的长边，向右的短边转为向上的短边。新车库位于C左上方，宽为3格、高为2格，标“车库”。旋转不会改变原来两条边的实际长度。
\BookFigure{figures/p072_s06-02.tex}
\end{BookAnswer}

% p072_s07
\BookQuestion{六、添一格的3种方法}
在下面的方格上再涂一个小方格，使涂色部分成为一个轴对称图形。你有几种涂法？请在下图中涂一涂，并画出它的对称轴。三幅图的已涂部分相同。\par
\BookFigure{figures/p072_s07-01.tex}
\begin{BookAnswer}
以下均从每幅网格左上角起，列从左向右、行从上向下计数。\par 第1种：涂第4行第3列，使底部横排向左补一格；对称轴是穿过第4列格子中心的竖直线。\par 第2种：涂第2行第5列；上下两排互为镜像，对称轴是穿过第3行中心的水平线。\par 第3种：涂第4行第6列，使两条直角臂各有3格；对称轴沿左下到右上的斜线，穿过第4行第4列和第2行第6列的中心。\par 逐一检验反射后每个已涂格都有对应格，得到这3种不同补法。
\BookFigure{figures/p072_s07-02.tex}
\end{BookAnswer}
